\documentclass[11pt]{article}
\usepackage{color}

%----------Packages----------
\usepackage{amsmath}
\usepackage{amssymb}
\usepackage{amsthm}
\usepackage{dsfont}
\usepackage{mathrsfs}
\usepackage{stmaryrd}
\usepackage[all]{xy}
\usepackage[mathcal]{eucal}
\usepackage{verbatim}  %%includes comment environment
\usepackage{fullpage}  %%smaller margins
\usepackage{multicol}
\usepackage{hyperref}
\usepackage{graphicx}

%Packages added to get the code block
\usepackage{listings}
\usepackage{xcolor}
\usepackage{tikz-cd}

\definecolor{codegreen}{rgb}{0,0.6,0}
\definecolor{codegray}{rgb}{0.5,0.5,0.5}
\definecolor{codepurple}{rgb}{0.58,0,0.82}
\definecolor{backcolour}{rgb}{0.95,0.95,0.92}

\lstdefinestyle{mystyle}{
    backgroundcolor=\color{backcolour},   
    commentstyle=\color{codegreen},
    keywordstyle=\color{magenta},
    numberstyle=\tiny\color{codegray},
    stringstyle=\color{codepurple},
    basicstyle=\ttfamily\footnotesize,
    breakatwhitespace=false,         
    breaklines=true,                 
    captionpos=b,                    
    keepspaces=true,                 
    numbers=left,                    
    numbersep=5pt,                  
    showspaces=false,                
    showstringspaces=false,
    showtabs=false,                  
    tabsize=2
}

\lstset{style=mystyle}


%---------Commands----------
\def\ZZ{\mathbb{Z}}
\def\QQ{\mathbb{Q}}
\def\PP{\mathbb{P}}
\def\NN{\mathbb{N}}
\def\OO{\mathcal{O}}
\def\AA{\mathbb{A}}
\def\BB{\mathbb{B}}
\def\FF{\mathbb{F}}
\def\LL{\mathbb{L}}
\def\RR{\mathbb{R}}
\def\CC{\mathbb{C}}
\def\GG{\mathbb{G}}
\def\lcm{\mathrm{lcm}}
\def\ord{\mathrm{ord}}
\def\rep{\mathrm{Rep}}
\def\dom{\mathrm{Dom}}
\def\ve{\varepsilon}
\def\gal{\mathrm{Gal}}
\def\emb{\mathrm{Emb}}

\newcommand{\dlegendre}[2]{\ensuremath{\left( \frac{#1}{#2} \right) }}
\newcommand{\tlegendre}[2]{\ensuremath{\big( \frac{#1}{#2} \big) }}

%----------Theorems----------
\theoremstyle{definition}
\newtheorem{theorem}{Theorem}[section]
\newtheorem{proposition}[theorem]{Proposition}
\newtheorem{lemma}[theorem]{Lemma}
\newtheorem{corollary}[theorem]{Corollary}
\newtheorem{definition}[theorem]{Definition}
\newtheorem{example}[theorem]{Example}
\newtheorem*{definition*}{Definition}
\newtheorem{nondefinition}[theorem]{Non-Definition}
\newtheorem{problem}{Problem}[section]
\newtheorem*{tip}{Pro Tip}
\newtheorem*{remark}{Remark}
\numberwithin{equation}{subsection}

\newenvironment{solution}
{%
  \vspace{0.5em}%
  \hrule%
  \vspace{0.2em}%
  \noindent\textbf{Solution:}\!\!\!\!\!
  \pushQED{\qed}
}
{%
  \popQED
  \vspace{.2em}
  \hrule%
  \vspace{0.75em}%
}

\newcommand{\separate}{
    \vspace{.2in}
    \hrule
    \vspace{.2in}
}


\newcommand{\head}[1]{
	\begin{center}
		{\large #1}
		\vspace{.05 in}
	\end{center}
	\bigskip 
}

\newcommand{\blue}[1]{\textcolor{blue}{#1}}
\newcommand{\red}[1]{\textcolor{red}{#1}}

%Problem Set Data
%number = 8
%course = 261
%semester = Fall 2025
%path = /Users/thyde/Documents/Teaching/Vassar/Semesters/2025 Fall/Math 261 Fall 2025/Problem Sets/Problem Set 8

\begin{document}
\pagestyle{plain}

%%---  sheet number for theorem counter
\setcounter{section}{3}   

\head{MATH 999, Summer 2026\\ Note 3: TBD\\
Trevor Hyde}

%HEADER

\begin{center}
   TBD
\end{center}
\vspace{.1in}

\begin{problem}
    Suppose that $L/K$ is a finite extension such that $|\gal(L/K)| = [L : K]$.
    We want to show that $L/K$ is Galois, which is to say that $L/K$ is the splitting extension of a separable polynomial.
    First note that $L = K(\alpha_1,\ldots, \alpha_n)$ for some $\alpha_i \in L$.
    Let $G := \gal(L/K)$.
    For each $\alpha_i$, let $H_i := \{\sigma \in G: \sigma(\alpha_i) = \alpha_i\}$.
    \begin{enumerate}
        \item \blue{Prove that $H_i$ is a subgroup of $G$.}
        \item Let 
        \[
            f_i(x) = \prod_{\sigma H_i \in G/H_i}(x - \sigma(\alpha_i)),
        \]
        where $G/H_i = \{\sigma H_i : \sigma\in G\}$ is the set of left cosets of $H_i$.
        \blue{Prove that $f_i(x)$ is separable for each $i$.}
        \item \blue{Prove that $f_i(x) \in K[x]$.}
        \textbf{Hint: }We need to know that $K$ is the fixed field of $\gal(L/K)$ for this. We don't know that yet.
        \item \blue{Prove that for $i \neq j$ we either have $f_i(x) = f_j(x)$ or $f_i(x)$ and $f_j(x)$ are coprime.}
        \item After possibly re-indexing, suppose that $f_1, f_2,\ldots, f_m$ are coprime and every $\alpha_i$ is a root of some $f_j$ with $1 \leq j \leq m$.
        Let $f(x) = f_1(x)\cdots f_m(x)$.
        \blue{Prove that }
    \end{enumerate}
\end{problem}


%Prove that for Galois extensions, the size of the group equals the degree of the extension.

%Examples


\end{document}